% Appendix
\label{sec:appendix}

\section*{Appendix A: Jacobian Derivation}\label{app:jacobian}

At an interior fixed point $v^* = (\rho_{\text{eff}}^*, \theta^{**}, N_{\text{eff}}^*)$ with $\rho_{\text{eff}}^* < \rho_1^*$, we define five partial derivatives evaluated at the fixed point. These coefficients govern the strength of each link in the feedback loop.

The first two coefficients capture how the Channel 2 information fraction responds to the other state variables:
\begin{align}
  a &\equiv \frac{d\mu_I^*}{d\theta^*} = -\frac{c_P}{(1-\theta^*)^2 \cdot |dF/d\mu_I|} < 0, \label{eq:coeff-a}\\
  b &\equiv \frac{d\mu_I^*}{d\rho_{\text{eff}}}\bigg|_{\theta^* \text{ fixed}} > 0. \label{eq:coeff-b}
\end{align}
The coefficient $a$ in~\eqref{eq:coeff-a} is negative because a higher crisis threshold raises the effective cost of private research, driving agents toward AI signals. The coefficient $b$ in~\eqref{eq:coeff-b} is positive because higher effective correlation reduces AI trading rents, making private information relatively more attractive.

The remaining three coefficients capture the Channel 3 liquidity link, the Channel 1 crisis sensitivity, and the $g_1$ amplification:
\begin{align}
  m &\equiv \frac{dg_3}{d\mu_I} = \frac{2N(N-1)\rho(1-\mu_I)}{(1+(N-1)\rho(1-\mu_I)^2)^2} > 0, \label{eq:coeff-m}\\
  h &\equiv \frac{d\theta^*}{d\rho_{\text{eff}}} > 0 \quad \text{(in the fragile region } \rho_{\text{eff}} \in (\rho_0, \rho_1^*)\text{)}, \label{eq:coeff-h}\\
  w &\equiv \frac{1-\rho}{N} > 0. \label{eq:coeff-w}
\end{align}
The coefficient $m$ in~\eqref{eq:coeff-m} measures how the effective market-maker count responds to changes in private information acquisition: fewer privately informed agents raise market-maker reliance on the common AI signal, reducing $N_{\text{eff}}$. The coefficient $h$ in~\eqref{eq:coeff-h} is the Channel 1 crisis sensitivity, which diverges as $\rho_{\text{eff}}$ approaches $\rho_1^*$. The coefficient $w$ in~\eqref{eq:coeff-w} captures the strength of the $g_1$ mapping from $N_{\text{eff}}$ to $\rho_{\text{eff}}$.

The Jacobian matrix of $T$ at $v^*$ is
\begin{equation}\label{eq:jacobian-full}
  DT = \begin{pmatrix} -wmb & -wm|a| & 0 \\ -hwmb & hwm|a| & 0 \\ mb & -m|a| & 0 \end{pmatrix},
\end{equation}
where the third column is zero because the updated state vector $(T_1(v), T_2(v), T_3(v))$ depends on the current $(\rho_{\text{eff}}, \theta^*)$ but not on the current $N_{\text{eff}}$ (which enters only through $g_1$, with $N_{\text{eff}}'$ computed from the current $\theta^*$ and $\rho_{\text{eff}}$). The structure of~\eqref{eq:jacobian-full} reveals that the system has rank two: one eigenvalue is zero, and stability depends entirely on the non-trivial eigenvalue $\ell_1 = w \cdot m \cdot (h|a| - b)$, as stated in~\eqref{eq:eigenvalue} of the main text.

\section*{Appendix B: Saddle-Node Bifurcation Verification}\label{app:saddle-node}

The Goldstein-Pauzner threshold equilibrium satisfies two conditions: an aggregate withdrawal equation $F_1 = 0$ and an indifference equation $F_2 = 0$, each in the unknowns $(x^*, \theta^*; \rho_{\text{eff}})$. Define the $2\times 2$ Jacobian
\begin{equation}\label{eq:h-divergence}
  J \;=\; D_{(x^*,\theta^*)}F \;=\; \begin{pmatrix} \partial F_1/\partial x^* & \partial F_1/\partial \theta^* \\ \partial F_2/\partial x^* & \partial F_2/\partial \theta^* \end{pmatrix}.
\end{equation}
The Morris-Shin uniqueness condition is equivalent to $\det(J) \neq 0$. At $\rho_{\text{eff}} = \rho_1^*$, the uniqueness condition $\alpha_{\text{SC}}\sqrt{\rho_{\text{eff}}/(1-\rho_{\text{eff}})} = 1$ binds, so $\det(J) = 0$: the unique threshold equilibrium and an emerging second equilibrium collide in a saddle-node bifurcation. We verify the three non-degeneracy conditions of the saddle-node normal form (\emph{cf.}\ Guckenheimer and Holmes, 1983, Theorem~3.4.1):
(a)~$\det(J) = 0$ at the bifurcation, by the Morris-Shin condition binding;
(b)~$\partial(\det J)/\partial \rho_{\text{eff}} \neq 0$, because $\det(J)$ is a smooth, strictly monotone function of $\alpha_{\text{SC}}^2\rho_{\text{eff}}/(1-\rho_{\text{eff}})$ near the critical value;
(c)~the quadratic coefficient in the reduced scalar equation is non-zero, which holds because $\theta^* = 1/2$ at the bifurcation point (since $\Phi(0) = 1/2$) and the second derivative $\partial^2 G/\partial(\theta^*)^2 \neq 0$ where $G$ is the scalar reduction obtained by eliminating $x^*$ via the implicit function theorem applied to $F_1$.

By the implicit function theorem applied to $F_1$, we can locally express $x^* = x^*(\theta^*, \rho_{\text{eff}})$ for $\rho_{\text{eff}} < \rho_1^*$, reducing the system to a single scalar equation $G(\theta^*, \rho_{\text{eff}}) \equiv F_2(x^*(\theta^*, \rho_{\text{eff}}), \theta^*; \rho_{\text{eff}}) = 0$. The saddle-node conditions~(a)--(c) imply the standard local expansion (Guckenheimer and Holmes, 1983, Theorem~3.4.1):
\begin{equation}\label{eq:saddle-expansion}
  \theta^*(\rho_{\text{eff}}) = \theta^*(\rho_1^*) - C\sqrt{\rho_1^* - \rho_{\text{eff}}} + O(\rho_1^* - \rho_{\text{eff}}), \quad C > 0,
\end{equation}
so that $h = d\theta^*/d\rho_{\text{eff}} = C/(2\sqrt{\rho_1^* - \rho_{\text{eff}}}) + O(1) \to \infty$ as $\rho_{\text{eff}} \to \rho_1^{*-}$. The local expansion~\eqref{eq:saddle-expansion} is the key asymptotic result: it establishes that $h$ diverges at the square-root rate characteristic of saddle-node bifurcations, which drives the amplification result in Proposition~\ref{prop:bifurcation}.
